\documentclass[11pt]{article}
\usepackage[a4paper,margin=1in]{geometry}
\usepackage{amsmath,amssymb,amsthm,mathtools}
\usepackage{graphicx}
\usepackage{hyperref}
\usepackage{microtype}

\newtheorem{theorem}{Theorem}
\newtheorem{lemma}{Lemma}
\newcommand{\Hc}{\mathcal H}
\newcommand{\B}{\mathcal B}
\newcommand{\Ran}{\operatorname{Ran}}

\title{Cline's formula fails for one-sided Drazin inverses}
\author{Run \texttt{fa\_banach\_001}, agent lane 10}
\date{2026-08-12}

\begin{document}
\maketitle

\section*{Source and answer}

Yan \cite{yan} proves that the nonzero parts of the one-sided Drazin spectra
of $ac$ and $ca$ agree, and asks in Question 3.10 whether the missing zero
point also agrees---equivalently, whether Cline's formula holds for
one-sided Drazin inverses in rings and one-sided generalized Drazin inverses
in Banach algebras.

\begin{center}
\includegraphics[width=0.96\linewidth]{figures/open_problem_crop.png}
\end{center}

The answer is no, simultaneously for the left and right versions and for
both Drazin and generalized Drazin invertibility.  The counterexamples lie
in the unital Banach algebra $\B(\ell^2)$, so they also answer the ring
version.

\section*{Definitions}

An element $T$ of a unital Banach algebra is \emph{left generalized Drazin
invertible} if some $X$ satisfies
\[
 TXT=XT^2,\qquad X^2T=X,\qquad TXT-T\ \hbox{is quasinilpotent}.
\]
Replacing ``quasinilpotent'' by ``nilpotent'' gives left Drazin
invertibility.  The right definitions reverse the displayed identities.
Every left invertible element is left Drazin invertible (of index zero), and
similarly on the right.

For $T\in\B(\Hc)$, write
\[
 H_0(T)=\left\{x\in\Hc:
   \lim_{m\to\infty}\|T^m x\|^{1/m}=0\right\}.
\]

\section*{Proof intuition}

Split the coordinates of $\ell^2$ into even and odd positions.  Let $D$
embed the whole space isometrically into the even coordinates.  Let $C$ be
the identity on those coordinates but shrink the $n$th odd coordinate by
$1/n$.  In the order $CD$, the shrinking part is never seen, so $CD=D$ is
left invertible.  In the reverse order $DC$, every vector first experiences
$C$ and is then pushed repeatedly through isometric even-coordinate
embeddings.  Thus no nonzero vector decays quasinilpotently, although the
odd basis vectors are sent to vectors of norm $1/n$.  A short idempotent
argument shows that a left generalized Drazin inverse would then have to be
an actual left inverse, which is impossible.

\section*{An algebraic obstruction}

\begin{lemma}\label{lem:obstruction}
If $T\in\B(\Hc)$ is left generalized Drazin invertible and $H_0(T)=\{0\}$,
then $T$ is left invertible.
\end{lemma}

\begin{proof}
Let $X$ be a left generalized Drazin inverse and put $Q=XT$.  The defining
identities give
\[
 Q^2=XTXT=X^2T^2=XT=Q,
 \qquad TQ=TXT=XT^2=QT.
\]
Moreover, $T-TQ=T(1-Q)$ is quasinilpotent.  Since $T$ commutes with $Q$,
for every $z\in\Ran(1-Q)$ and every $m\ge1$,
\[
 T^m z=[T(1-Q)]^m z.
\]
The spectral-radius formula therefore implies
$\Ran(1-Q)\subseteq H_0(T)=\{0\}$.  Hence $Q=1$, so $XT=1$ and $T$ is left
invertible.
\end{proof}

\section*{The counterexample}

\begin{theorem}\label{thm:main}
There exist $a,c\in\B(\ell^2(\mathbb N))$ such that $ac$ is left Drazin
invertible but $ca$ is not left generalized Drazin invertible.  There also
exist such a pair with ``left'' replaced everywhere by ``right''.
Consequently, Cline's formula fails for every one-sided class asked about in
Question 3.10.
\end{theorem}

\begin{proof}
Let $(e_n)_{n\ge1}$ be the standard basis of $\Hc=\ell^2(\mathbb N)$ and
define $D,C\in\B(\Hc)$ by
\[
 De_n=e_{2n},\qquad
 Ce_{2n}=e_{2n},\qquad
 Ce_{2n-1}=\frac1n e_{2n-1}\quad(n\ge1).
\]
The operator $D$ is an isometry, $D^*D=1$, and $C$ acts as the identity on
$\Ran D$.  Thus
\[
 CD=D.
\]
In particular, $CD$ is left invertible, with left inverse $D^*$, and hence
is left Drazin invertible.

Set $T=DC$.  From $CD=D$, induction gives
\[
 T^m=(DC)^m=D^mC\qquad(m\ge1).
\]
Because every $D^m$ is an isometry,
\[
 \|T^m x\|=\|Cx\|\qquad(x\in\Hc,\ m\ge1).
\]
The diagonal operator $C$ is injective.  Hence for $x\ne0$,
\[
 \lim_{m\to\infty}\|T^m x\|^{1/m}
 =\lim_{m\to\infty}\|Cx\|^{1/m}=1,
\]
and therefore $H_0(T)=\{0\}$.

On the other hand,
\[
 \|Te_{2n-1}\|=\frac1n\longrightarrow0.
\]
Thus $T$ is not bounded below.  It cannot have a bounded left inverse,
because $XT=1$ would imply
$\|x\|\le\|X\|\,\|Tx\|$ for every $x$.  Lemma \ref{lem:obstruction} now
shows that $T=DC$ is not left generalized Drazin invertible.  Taking
$a=C$ and $c=D$ proves the left assertion, and also disproves the weaker
left-Drazin closure.

For the right assertion, take $a=D^*$ and $c=C=C^*$.  Then
\[
 ac=D^*C=(CD)^*=D^*,
 \qquad ca=CD^*=(DC)^*.
\]
The coisometry $D^*$ is right invertible, with right inverse $D$.  If
$(DC)^*$ were right generalized Drazin invertible, taking adjoints in its
three defining identities would make $DC$ left generalized Drazin
invertible; adjoints preserve quasinilpotence in $\B(\Hc)$.  This
contradicts the first half of the proof.  Hence the right versions fail as
well.
\end{proof}

\section*{Later-literature audit and scope}

A later paper of Kolund\v{z}ija and Mosi\'c \cite{km}, published online on
21 October 2025, studies conditional Cline formulas for the same one-sided
generalized inverses.  Its Example 3.1 is intended to show that a commutation
condition cannot be omitted and asserts a left generalized-Drazin failure.
That particular example does not establish the assertion.  In its notation,
the reversed shift satisfies
\[
 BAe_1=0,\qquad BAe_n=e_{n+1}\quad(n\ge2).
\]
Thus the reducing decomposition
\[
 \ell^2=\operatorname{span}\{e_1\}
 \oplus \overline{\operatorname{span}}\{e_n:n\ge2\}
\]
writes $BA$ as $0\oplus S$, where $S$ is a unilateral isometry.  Hence
$BA$ is left Drazin invertible.  The other diagonal block in their displayed
operator is quasinilpotent, hence left generalized Drazin invertible with
inverse zero; their whole reversed product is therefore left generalized
Drazin invertible.  A separate proof-gap record accompanies this packet.

The construction in Theorem \ref{thm:main} repairs that obstruction and is
self-contained.  Searches through 2026-08-12 of the four run indexes, the
exact question and title, citations, and targeted one-sided Cline-formula
literature found no other valid unconditional counterexample.  This is a
bounded novelty check, not an exhaustive survey of all non-indexed
literature.

The theorem resolves the existence/equivalence part of Question 3.10 in the
negative for both orientations and both notions.  It does not characterize
the additional hypotheses under which a one-sided Cline formula is valid.

\section*{Human review recommendation}

\enlargethispage{4\baselineskip}
Review as a likely-valid full counterexample.  The decisive checks are the
identity $(DC)^m=D^mC$, the conclusion $H_0(DC)=\{0\}$, and Lemma
\ref{lem:obstruction}.  The later-paper proof-gap note should be reviewed
independently from the validity of the new construction.

\begin{thebibliography}{9}
\bibitem{yan}
Kai Yan,
\emph{On the spectral identities and fundamental properties of one-sided
Drazin inverses in Banach algebras}, arXiv:2504.18995 (2025), Question 3.10.

\bibitem{km}
Milica Z. Kolund\v{z}ija and Dijana Mosi\'c,
\emph{Properties of one-sided generalized Drazin inverses in Banach
algebras}, Advances in Operator Theory \textbf{11}, Article 2 (2026),
published online 21 October 2025,
\href{https://doi.org/10.1007/s43036-025-00479-1}{doi:10.1007/s43036-025-00479-1}.
\end{thebibliography}

\end{document}
